\documentclass[11pt]{article}

\usepackage[margin=1in]{geometry}
\usepackage{array}
\usepackage{booktabs}
\usepackage{enumitem}
\usepackage{hyperref}
\usepackage{listings}
\usepackage{microtype}
\usepackage{needspace}
\usepackage{tabularx}
\usepackage[T1]{fontenc}
\usepackage{lmodern}
\usepackage{xcolor}

\hypersetup{
    colorlinks=true,
    linkcolor=blue,
    urlcolor=blue
}

\definecolor{codebg}{RGB}{247,247,247}
\definecolor{codekw}{RGB}{45,76,150}
\definecolor{codecomment}{RGB}{80,115,80}
\definecolor{codestring}{RGB}{145,70,45}
\definecolor{decisionbg}{RGB}{240,245,250}

\newcolumntype{Y}{>{\raggedright\arraybackslash}X}
\newcolumntype{L}[1]{>{\raggedright\arraybackslash}p{#1}}

\lstdefinelanguage{Rust}{
    morekeywords={as,async,await,break,const,continue,crate,dyn,else,enum,extern,
        false,fn,for,if,impl,in,let,loop,match,mod,move,mut,pub,ref,return,
        self,Self,static,struct,super,trait,true,type,unsafe,use,where,while},
    sensitive=true,
    morecomment=[l]{//},
    morecomment=[s]{/*}{*/},
    morestring=[b]"
}

\lstset{
    language=Rust,
    basicstyle=\ttfamily\small,
    keywordstyle=\color{codekw}\bfseries,
    commentstyle=\color{codecomment}\itshape,
    stringstyle=\color{codestring},
    backgroundcolor=\color{codebg},
    frame=single,
    framerule=0.3pt,
    rulecolor=\color{black!25},
    columns=fullflexible,
    keepspaces=true,
    showstringspaces=false,
    breaklines=true,
    tabsize=4,
    xleftmargin=0.5em,
    xrightmargin=0.5em,
    aboveskip=0.8em,
    belowskip=0.8em
}

\title{\textit{cs-mail} Desktop MVP and CLI Administration\\
\large AI-Agent Milestone and Build Plan}
\author{C-SQD}
\date{Draft 0.2 --- September 2026}

\setlength{\emergencystretch}{1.5em}

\begin{document}
\maketitle

\begin{abstract}
This plan develops a native desktop pilot, an operator administration CLI, and a
reusable client core around the revised relationship-request model. Alice makes
one simulated charge for a request to Bob, sends permitted follow-ups without
another charge, and sees the outcome of his decision. The service then carries
a simulated forfeiture through maturity, quarterly allocation, and payment
recovery. Each milestone requires a usable journey and evidence of failure
recovery. The pilot introduces no real redeemable value.
\end{abstract}

\section{Start with the Revised Contract}

The \textit{Protocol Specification} governs requests and permission. The normative
C-SQD financial-program section of the \textit{Deployment and Migration Profile}
governs the pool, corporate allocation, and equal member distributions. The
\textit{Rust Reference Architecture} explains the intended implementation shape.
This plan is nonnormative and does not claim that the existing Rust code already
implements the new contracts.

Before product milestones begin, establish the revised semantic baseline:
version the request and obligation contracts, document treatment of any old
quotes or records, and replace superseded financial expectations with explicit
conformance vectors. Preserve applicable security, privacy, ordering, and
conservation guarantees. An old test requiring separate bonds for follow-up
messages must not override the revised specification merely because it predates
the change.

The baseline must model request-specific payment operations, refunds against the
original charge, pending forfeitures, restricted member funds, and finalized
payables. One pending request permits at most one charge; recipient settings
determine whether it admits more messages. The base draft uses no nonzero
persistence reserve. Proposed rates, eligibility rules, and time periods use
explicit simulation-policy versions and must not be presented as approved live
terms.

Once that baseline is established, product work extends it through small vertical
slices. Further changes to value semantics, signatures, identity meaning, or
privacy scope require a decision record explaining the user-visible need,
affected contracts, compatibility, migration, and validation. Routine reversible
implementation choices remain the agent's responsibility.

\section{The Pilot We Are Building}

The initial audience is an invite-only group of approximately 20--50 cooperative
users. The product is a macOS desktop application plus an administration CLI,
served by one provider, one region, and one PostgreSQL deployment. Windows follows
after the macOS journey is stable. Later Swift and Kotlin clients reuse a
platform-neutral Rust client core through a narrow binding.

The value model is a deterministic external-payment simulator. It can approve,
decline, delay, or report an uncertain request charge; confirm or fail a refund;
and execute or fail a finalized member payment. All amounts are explicitly
nonredeemable. Users do not receive a generally spendable test-credit wallet,
and the operator configures simulator scenarios rather than topping up accounts.

Native content remains encrypted at endpoints. Devices hold revocable key
authority, recovery uses user-held material or trusted-device transfer, and the
provider cannot decrypt content. The desktop stores ciphertext and minimal
metadata by default, with no persistent plaintext body or search index.

SMTP, independent-provider federation, public signup, live payments, a mobile
release, and multi-region operation are outside this pilot. The financial
simulation is included because it validates the revised business rules; it is
not evidence that production payment operations are ready.

\section{Keep Shared Mechanics and Product Journeys Separate}

The stable foundation remains organized around primitives, authentication,
ledger, request transitions, content, wire, storage, and workers. Product modules
coordinate a user journey without implementing another version of those rules.

\begin{lstlisting}[language={}]
crates/
  existing cs-mail-* crates       # revised semantic foundation
  cs-mail-api-types               # versioned transport contracts
  cs-mail-client-core             # portable client use cases
  cs-mail-mobile-ffi              # later portability milestone

apps/
  cs-mail-server                  # API library and cs-maild
  cs-mail-admin-cli               # authenticated operator client
  cs-mail-desktop                 # Tauri shell and platform adapters
\end{lstlisting}

Dependencies point inward: applications use the shared core, which uses the
foundation; foundation crates never depend on applications. API types should not
depend on a web framework, desktop shell, concrete database, or operating-system
keychain. Create another crate only for a stable reusable contract, dependency
boundary, or independently built artifact, not for every new feature.

Within client core and server, organize identity/device management, requests,
messages, relationships, financial activity, distributions, recovery, and
administration as vertical modules. Request orchestration handles quote,
authorization, content upload, payment progress, admission, and retry. Financial
views consume canonical obligation status; they do not reconstruct settlement
from UI events. The member program consumes eligibility and funding evidence
separately from message handling.

The client core defines ports for key custody, local ciphertext storage,
transport, clock, and notifications. Tauri, Swift, Kotlin, Keychain, SQLite, and
concrete HTTP libraries belong in adapters. Use constrained, versioned DTOs and
opaque canonical signed bytes for protocol mutations. An HTTP handler cannot
translate a signed command into a different meaning.

\section{Contracts That Apply to Every Milestone}

Endpoint private keys use operating-system secure storage and never enter logs,
configuration, SQLite, provider APIs, or diagnostics. Device enrollment and
revocation require authenticated authority. Administrative credentials are distinct
from user credentials, and the CLI invokes an authenticated API rather than
writing canonical tables directly.

Every visible outcome has a durable cause and stable retry identity. Unknown
payment or commit outcomes remain distinguishable from failure. UI language
separates charge preparation, an open request, a refund due, refund completion,
finalized member allocation, and payment completion. No screen suggests that
rejecting a sender pays the recipient or that a payable can fund another request.

Use separate protocol, wire, API, event, policy, and database versions. Migrations
are forward-only after merge. Logs contain safe scoped references and typed
errors, not plaintext, secrets, identity mappings, full signed commands, or
complete financial snapshots. Retention policy must be explicit before pilot
data is created, including records needed for unsettled simulated obligations.

At each implementation handoff, run the workspace quality gate:
\begin{lstlisting}[language={}]
cargo fmt --all -- --check
cargo clippy --workspace --all-targets -- -D warnings
cargo test --workspace
\end{lstlisting}
PostgreSQL integration tests run against an isolated database; ignored tests do
not count as evidence. Add contract, migration, desktop smoke, and recovery tests
as the relevant surfaces appear.

\section{Milestone Map}

\begin{center}
\small
\begin{tabularx}{\textwidth}{@{}L{0.09\textwidth}L{0.29\textwidth}Y@{}}
\toprule
Stage & User-visible result & Exit evidence \\
\midrule
Baseline & Revised request and financial contracts & Versioned decisions and conformance vectors \\
0 & Product foundation is ready & Buildable skeleton and dependency checks \\
1 & Operator starts the service & Bootstrap, simulator, CLI, and restart demo \\
2 & Alice owns an identity & Enrollment, key custody, and reopen demo \\
3 & Alice requests a relationship & One simulated charge and encrypted admission \\
4 & Bob controls the request & Decisions, follow-up policy, and accurate activity views \\
5 & Time and failure are recoverable & Deadline, provider-event, and crash matrix \\
5b & The member cycle is demonstrable & Reproducible quarter and failed-payment recovery \\
6 & A small pilot can use it & Signed builds, recovery, restore, and usability evidence \\
7 & The core is portable & Swift/Kotlin binding parity harness \\
\bottomrule
\end{tabularx}
\end{center}

Each stage depends on the preceding one. A milestone is complete when its
demonstration runs from a clean checkout and its relevant failure cases pass,
not merely when the new code compiles.

\section{Milestone 0: Establish the Product Foundation}

After the semantic baseline, create the product packages and enforce inward
dependencies. Define shared package metadata, lints, and dependency versions.
Document the architecture boundary and add a command that runs the full quality
gate without hiding individual checks. Empty server and CLI binaries should
report version and help.

Demonstrate the workspace build and dependency graph from a clean checkout.
The exit evidence is a passing baseline suite, reviewed dependency rules, and a
small package skeleton. No enrollment, API listener, desktop workflow, or further
semantic change is part of this stage.

\section{Milestone 1: Bring the Service and Simulator to Life}

The operator can migrate an empty database, bootstrap one provider idempotently,
start the API, create invites, inspect policy, and select simulation scenarios
without editing SQL. Use Tokio, Axum, Rustls, structured tracing, and Clap as the
initial process, HTTP, TLS, diagnostics, and CLI choices. Keep synchronous storage
behind a bounded blocking boundary until a tested replacement is appropriate.

Configuration identifies the deployment domain, provider authority, public URL,
database, content limits, retention versions, and simulation mode. Ordinary
administration uses an authenticated API. A local bootstrap/migration command
may initialize the deployment, but cannot silently replace an existing identity
or key registry.

The simulator persists operation identities and supports controlled delayed,
duplicate, out-of-order, and failed results. Its scenario changes are audited
and separate from canonical financial writes. No live provider credentials or
redeemable amounts are enabled.

Tests cover migration from empty, repeated bootstrap, missing keys, user/admin
authority separation, invite replay, simulator-operation idempotency, and readiness
that distinguishes process health from database readiness. Demonstrate restart
without duplicate authority or simulated operations. Exit evidence includes a
clean-start transcript, API contracts, database tests, and redacted diagnostics.

\section{Milestone 2: Give Alice an Identity She Controls}

Alice installs the macOS application, redeems an invite, creates a protocol
identity and device keys, and reopens the same identity after quitting. Client
use cases coordinate identity resolution, enrollment, device listing, and
revocation through portable ports. The Tauri shell remains thin.

Store only key references, ciphertext, cursors, minimized relationship metadata,
and nonsecret preferences in SQLite. Use the macOS keychain for secrets. Explain
the recovery capabilities that are available at this stage without implying
provider access to old content.

Tests reject foreign-scope or revoked device authority, conflicting invite reuse,
and secret leakage to storage or logs. Local schema upgrades preserve ciphertext
and cursors. Demonstrate enrollment, keychain denial, quit/reopen, and device
listing on a clean user account. Sending and completed recovery remain deferred.

\section{Milestone 3: Submit One Encrypted Relationship Request}

Alice resolves Bob and reviews $C$, $S$, the total conditional charge, the request
deadline rule, and possible refunds. The UI identifies the value as simulated.
She authorizes one request, encrypts the initial message, and obtains signed
evidence of admission after the simulator confirms capture and delivery
preparation is durable.

The client state machine persists the request and operation references across
quote, create, payment preparation, ciphertext upload, and admission. An existing
preparing or open request is resumed instead of charged again. No delivery is
visible before durable admission. Bob receives an encrypted inbox record and
decrypts at his endpoint.

Test two-device creation with distinct command keys, capture failure, uncertainty,
expired terms, wrong scope, modified declarations, ciphertext tampering, and
network loss at each boundary. A declined charge admits nothing; an unknown
charge is reconciled under its original identity. Alice's request activity must
match the ledger and receipt, without a spendable balance view.

Demonstrate the complete encrypted path and recovery after an interrupted charge
response. Exit evidence is an end-to-end request test, a payment-step retry
matrix, and plaintext-absence checks. Attachments, search, rich text, SMTP, and
real currency remain deferred.

\section{Milestone 4: Make Bob's Decision and Follow-Up Controls Real}

Bob sees one request with its grouped messages and a clear decision deadline.
Acceptance grants future communication and records one full refund obligation;
rejection records processing revenue and pending forfeiture; blocking additionally
prevents contact. The UI explains permission first, then money, and never
describes rejection as a payment to Bob.

Implement recipient follow-up preferences with bounded count/rate controls and
canonical admission enforcement. Alice can add a message only when the identified
request and current policy allow it. Its own delivery state remains visible, but
the charge and common deadline remain unchanged. Closing or dismissing the
recipient's prompt is not a decision.

Show request charges and refund status separately from member distribution
statements. Accepted sending uses standing permission rather than a zero-value
quote. Block/unblock management preserves the difference between rejection and
prohibition, and projection rebuild after restart reproduces canonical state.

Tests cover one request with three messages and one refund, disabled follow-ups,
allowance replay, policy-change races, decision races, alias history, and block
before versus after admission. A follow-up refusal must not start another charge.
Demonstrate accept, reject, block/unblock, and allowed/disabled follow-ups from
both clients. Exit evidence includes reviewed financial wording and encrypted
accepted-message delivery.

\section{Milestone 5: Recover Across Time and Failure}

Run scheduled request deadlines, outbox delivery, payment reconciliation, and
retention through supervised leased workers. Jobs recheck current state and use
stable operation identities. Configure a distinct scheduler authority, bounded
retries, graceful shutdown, database TLS, and a bounded connection pool without
changing the revised semantics.

This milestone makes uncertainty visible. Kill the service after capture but
before its response is recorded; cancel a preparing request before delayed capture
confirmation; replay provider notifications out of order; and fail a refund before
retrying it. The request stays cancelled when appropriate, all captured money is
returned once, and a terminal obligation does not disappear during recovery.

Exercise timely acceptance delayed behind expiry processing and late acceptance
after a closed window. Preserve current permission while materializing old
financial outcomes. Offline recipients and cursor replay must not duplicate
delivery or financial effects. Retention must protect content still needed by
undelivered work while eventually deleting eligible records.

Demonstrate restart, lease reclaim, deadline processing, delayed-provider recovery,
and one visible message/economic outcome after retries. Exit evidence is a
failure-injection matrix, scheduler-authority tests, a sustained concurrency run,
and an operator reconciliation procedure.

\section{Milestone 5b: Complete a Simulated Quarter}

The operator can demonstrate how a rejected request eventually benefits members
without paying its recipient directly. Create simulated forfeiture lots with
different maturity and quoted-rate versions, and a synthetic membership set with
explicit activity and verification evidence. These fixture policies are visibly
provisional; they do not settle live legal-entity eligibility or identity checks.

Prepare a quarter from fixed funding and eligibility manifests. Apply the
financial profile's once-only corporate calculation, equal minor-unit allocation,
and restricted remainder. Publish the finalized batch atomically and show an
aggregate reconciliation plus each member's own statement. Operator access must
not expose the correspondence behind another member's funding or eligibility.

Test refundable and immature exclusions, duplicate members and aliases, exact
cutoff boundaries, multiple rate cohorts, zero eligible members, empty funding,
rounding, and replay after partial preparation. The next quarter must not assess
previously assessed funds again. A delayed or below-threshold payable remains
assigned to its member and cannot become spendable request funding.

Finally fail one simulated payout, cross the quarter boundary, and recover it.
The member's fixed entitlement remains intact throughout. Exit evidence includes
reproducible manifests, conservation checks, a failed-payment demo, and a clear
distinction between unallocated carryforward and unpaid finalized obligations.
Live payouts and production verification remain deferred.

\section{Milestone 6: Run a Small Supported Pilot}

Package and sign the macOS application and versioned server/CLI artifacts.
Complete user-held recovery and trusted-device enrollment, with explicit device
revocation and no provider content-decryption authority. Add understandable
onboarding for requests, refunds, simulated distributions, follow-up preferences,
encryption, recovery, and device loss.

Provide scoped diagnostics export, support and abuse intake, deletion procedures,
backup and restore, incident escalation, and release rollback instructions.
Backups and retained obligation records must respect separate keys and access
roles. Windows adapters follow only after the core macOS acceptance path is
stable and platform-neutral.

Tests cover clean installation and upgrade, recovery on a new device, revoked
device refusal, diagnostics secret scans, and restore followed by reconciliation
of requests, refunds, quarter manifests, payables, and inbox state. Run a
50-identity soak with offline periods, decisions, retries, and restarts.

Give nondeveloper testers clean machines and the pilot guide. They should enroll,
request a relationship, use follow-up settings, decide, inspect financial status,
recover a device, and export safe diagnostics without direct database repair.
Exit evidence includes signed releases, restore and recovery records, and resolved
critical usability findings. Public signup and production financial value remain
outside this milestone.

\section{Milestone 7: Prove Mobile Portability}

Expose a narrow FFI facade for identity, request progress, decisions, encrypted
inbox, financial status, distribution statements, and device management. It uses
bounded values, stable error codes, cancellation, and explicit ownership rather
than exporting internal domain structs or secret bytes.

Generate Swift and Kotlin bindings and provide mock platform adapters for secure
keys, transport, ciphertext storage, and notifications. Run the same client-core
stories through Rust and each binding, including a pending refund and failed
member payment. Malformed input returns typed errors rather than panics, and
process termination cannot leave an unknown operation represented as completed.

The exit evidence is a versioned facade, clean binding builds, and behavior parity
without Tauri in the shared dependency graph. Store releases, mobile push, and a
polished mobile UI are separate work.

\section{How an Implementation Agent Should Proceed}

For each slice, read the applicable revised contracts and local instructions,
name one user-visible outcome, identify its trust and persistence boundaries,
then implement contracts, domain orchestration, adapters, and UI in that order.
Use a real end-to-end path early; do not build every route before demonstrating
one request. Report what changed, what was demonstrated, and what remains open.

Do not bypass signed execution, write canonical tables from the CLI, persist
plaintext for convenience, hide an uncertain charge, reinterpret old migrations,
or weaken a conformance test merely to finish a milestone. Changes to the approved
baseline require an explicit decision and corresponding new evidence. Scope must
not expand to live value or SMTP simply because an adapter is available.

If reconciliation fails, preserve the affected evidence and investigate the
specific obligation. If secrets leak, stop accumulating affected user data until
the boundary is repaired. If a policy parameter is unresolved, use an explicitly
versioned simulation fixture and keep the live deployment gate closed rather
than presenting that fixture as a business commitment.

\section{Completion and the First Assignment}

The pilot is complete when clean installations can enroll and recover identities,
exchange encrypted relationship requests and accepted messages, enforce follow-up
preferences, survive deadlines and failures, and reproduce a simulated quarter
with correct unpaid obligations. The operator can deploy, reconcile, restore,
and support it without arbitrary database edits. The published pilot policy
states simulation limits, retention, abuse review, and deletion behavior.

The first assignment is therefore to establish the revised semantic baseline
and its conformance vectors. After that, Milestone~0 and the smallest service/CLI
walking skeleton provide a bounded product start. Preserving the earlier
per-message bond model is not a prerequisite for this revised plan; preserving
the agreed guarantees under the new request model is.

\end{document}
